\begin{proof}[Proof of \cref{obj:thm:connected-interval-frontier}]
Fix \(\alpha\in(0,1)\). Throughout the proof, write
\[
\beta:=1-\alpha>0,\qquad
N:=nq>0,\qquad
\ell:=\log(e+N)\ge1,
\]
where \(n,d\ge1\), \(0<q\le1\), and \(q\ell^2\le1/64\).
Define the two length scales by
\[
s_{\mathrm{par}}:=\min\{1,N^{-1/2}\},
\qquad
s_{\mathrm{many}}:=\min\left\{1,\frac{d}{N\ell}\right\}.
\]

\begin{enumerate}
\item \label{proof:connected-interval-upper}\textbf{Upper bound for the risk-envelope interval.}
By \cref{obj:thm:mixed-count-upper}, there is a universal
\(\overline C>0\) such that, for every \(P\in\mathcal M_{n,d,q}\),
\[
E_P\!\left[
  \bigl(\widehat\tau^{\mathrm{mix}}_{n,d,q}-\tau(P)\bigr)^2
\right]
\le \mathfrak U_{n,d,q}
\le \overline C\,r_{n,d,q}.
\]
For every observed sample, define the envelope radius by
\[
\rho_{\mathrm{env}}
:=\sqrt{\frac{\mathfrak U_{n,d,q}}{\alpha}}.
\]
The endpoints in \cref{obj:def:risk-envelope-interval} satisfy
\[
\operatorname{length}(I^{\mathrm{mix}}_\alpha)
=
\min\{1,\widehat\tau^{\mathrm{mix}}_{n,d,q}+\rho_{\mathrm{env}}\}
-
\max\{-1,\widehat\tau^{\mathrm{mix}}_{n,d,q}-\rho_{\mathrm{env}}\}
\le 2\rho_{\mathrm{env}}.
\]
Consequently, with
\[
C_\alpha:=2\sqrt{\frac{\overline C}{\alpha}}>0,
\]
we have
\[
E_P[\operatorname{length}(I^{\mathrm{mix}}_\alpha)]
\le
2\sqrt{\frac{\mathfrak U_{n,d,q}}{\alpha}}
\le C_\alpha\sqrt{r_{n,d,q}}.
\]
Taking the supremum gives
\[
\sup_{P\in\mathcal M_{n,d,q}}
E_P[\operatorname{length}(I^{\mathrm{mix}}_\alpha)]
\le C_\alpha\sqrt{r_{n,d,q}}.
\]
For an empty model class, the supremum is zero by the convention in
\cref{obj:def:interval-length-risk}, and the same inequality follows from
\(r_{n,d,q}\ge0\).
% lean: CausalSmith.Stat.MarRareqLogfrontier.riskEnvelope_interval_worst_upper

\item \label{proof:connected-interval-calibration}\textbf{Calibration of the activated regime.}
Define constants depending only on \(\alpha\) by
\[
\eta:=e^{-32},\qquad
M_{\mathrm{act},\alpha}:=\left\lceil\frac8\beta\right\rceil+1,
\qquad
D_{\mathrm{act},\alpha}
:=
\left\lceil
\frac{18432M_{\mathrm{act},\alpha}^2}{\beta}
\right\rceil+1,
\]
\[
T_{\mathrm{act},\alpha}
:=
\max\left\{
1,e^2,4\eta,2\eta D_{\mathrm{act},\alpha},
\exp\left(\frac{32-\log(\beta/16)}{28}\right)
\right\},
\qquad
c_{\mathrm{act},\alpha}
:=\frac{21\beta}{32}\frac{\eta}{24}>0.
\]
The ceiling inequalities give
\[
M_{\mathrm{act},\alpha}\ge1,\qquad
\frac1{M_{\mathrm{act},\alpha}}\le\frac\beta8,\qquad
D_{\mathrm{act},\alpha}\ge1,\qquad
\frac{18432M_{\mathrm{act},\alpha}^2}{\beta}
\le D_{\mathrm{act},\alpha}.
\]
Use the quantities from \cref{obj:def:activation-lower-handle}:
\[
K:=\lceil\ell\rceil,\qquad H:=K^2,\qquad
b_0:=\frac{\eta}{N\ell},\qquad
J:=\min\left\{d-1,\left\lfloor\frac1{2b_0}\right\rfloor\right\}.
\]
They are defined for all the parameters under consideration, including
\(d=1\), when \(J=0\).

Suppose first that \(N\ge T_{\mathrm{act},\alpha}\). Since \(\ell\ge1\),
\[
D_{\mathrm{act},\alpha}
\le \frac{N}{2\eta}
\le \frac{N\ell}{2\eta}
=\frac1{2b_0}.
\]
Also,
\[
b_0\le\frac{\eta}{N}\le\frac14,\qquad
\ell\ge2,\qquad K\ge2.
\]
The ceiling bound \(K\le\ell+1\le2\ell\), together with the slice restriction,
gives
\[
qH=qK^2\le4q\ell^2\le\frac1{16}.
\]
By \cref{obj:lem:moment-matched-reciprocal-priors}, choose finitely supported
probability laws \(\Pi_0,\Pi_1\) on \([1,H]\) with matching moments through
degree \(K\). Interchange their names if necessary and define
\[
\mu_{\mathrm{rec},i}:=\int z^{-1}\,d\Pi_i(z)
\quad(i=0,1).
\]
Then
\[
\int z^v\,d\Pi_0(z)=\int z^v\,d\Pi_1(z)
\quad(0\le v\le K),
\qquad
\mu_{\mathrm{rec},1}-\mu_{\mathrm{rec},0}\ge\frac1{12}.
\]

We will show that whenever \(J\ge D_{\mathrm{act},\alpha}\),
\[
c_{\mathrm{act},\alpha}s_{\mathrm{many}}
\le\mathfrak L_{n,d,q,\alpha}.
\]
% lean: CausalSmith.Stat.MarRareqLogfrontier.interval_activated_large_regime

\item \textbf{The activated grid and its likelihood comparison.}
Continue under
\(N\ge T_{\mathrm{act},\alpha}\) and \(J\ge D_{\mathrm{act},\alpha}\).
For each integer \(0\le\iota\le M_{\mathrm{act},\alpha}\), define
\[
\Pi^{(\iota)}
:=
\left(1-\frac{\iota}{M_{\mathrm{act},\alpha}}\right)\Pi_0
+\frac{\iota}{M_{\mathrm{act},\alpha}}\Pi_1.
\]
These are finitely supported probability laws on \([1,H]\), with the same
moments through degree \(K\).

For completeness, the full-data laws underlying the activated experiment
can be specified as follows. For a fixed
\(\mathbf z=(z_0,\ldots,z_{d-1})\in[1,H]^d\), let
\(P^{\mathrm{act}}_{\mathbf z}\) have baseline probabilities
\[
P^{\mathrm{act}}_{\mathbf z}(X=x)
=
\begin{cases}
b_0,&0\le x<J,\\
1-Jb_0,&x=J,\\
0,&J<x<d.
\end{cases}
\]
Let treatment be an independent fair coin, and set
\[
S(0)=S(1)=Y(0)=0,\qquad S=0,\qquad Y=AY(1).
\]
Conditional on \(X=x\), generate the treated potential outcome and arrival
independently, with
\[
P^{\mathrm{act}}_{\mathbf z}(Y(1)=1\mid X=x)
=
\begin{cases}
z_x^{-1},&x<J,\\
0,&x\ge J,
\end{cases}
\qquad
P^{\mathrm{act}}_{\mathbf z}(R=1\mid X=x,A)
=
\begin{cases}
qz_x,&x<J,\\
qH,&x\ge J.
\end{cases}
\]
The inequalities \(Jb_0\le1/2\) and \(qH\le1/16\) ensure valid
probabilities. Independence of treatment, consistency, conditional
independence of arrival and outcome, and the arrival floor follow from
this construction. Thus
\[
P^{\mathrm{act}}_{\mathbf z}\in\mathcal M_{n,d,q},
\qquad
\tau(P^{\mathrm{act}}_{\mathbf z})
=b_0\sum_{x=0}^{J-1}z_x^{-1}.
\]

Let \(\widetilde{\mathbb Q}_{\mathrm{act},\iota}\) be the augmented
sample law obtained from the rule in
\cref{obj:def:activation-lower-handle}, using
\((\Pi^{(\iota)})^{\otimes d}\) for the latent vector. Let
\(\mathbb Q_{\mathrm{act},\iota}\) be its projection onto the observed
sample. The conditional probabilities in that construction give
\[
\mathbb Q_{\mathrm{act},\iota}
=
\int
\mathcal L_{P^{\mathrm{act}}_{\mathbf z}}(\mathbf O_n)\,
d(\Pi^{(\iota)})^{\otimes d}(\mathbf z).
\]
Define the activation count in each rare category and the common cutoff
event by
\[
V_x:=\sum_{i=1}^n
\mathbf1\{X_i=x,\ \mathsf B_{\mathrm{act},i}=1\}
\quad(0\le x<J),
\qquad
G_{\mathrm{act}}:=\{V_x\le K\text{ for every }x<J\}.
\]
Conditional on the membership, treatment, and activation pattern, every
activated observation in category \(x\) contributes an affine factor in
\(z_x\), and every inactive observation contributes a constant factor.
The conditional likelihood therefore has degree at most \(V_x\) in each
\(z_x\). On \(G_{\mathrm{act}}\), integration against the product priors
and moment matching imply equality of all augmented sample atoms:
\[
\widetilde{\mathbb Q}_{\mathrm{act},\iota}(\{\mathbf o_{\mathrm{aug}}\})
=
\widetilde{\mathbb Q}_{\mathrm{act},0}(\{\mathbf o_{\mathrm{aug}}\})
\qquad(\mathbf o_{\mathrm{aug}}\in G_{\mathrm{act}}).
\]
Hence the two laws also assign the same mass to \(G_{\mathrm{act}}\)
and its complement. Splitting any event into its intersections with
these two sets, and then projecting, yields
\[
\operatorname{TV}
\bigl(\mathbb Q_{\mathrm{act},0},
      \mathbb Q_{\mathrm{act},\iota}\bigr)
\le
\widetilde{\mathbb Q}_{\mathrm{act},0}(G_{\mathrm{act}}^c).
\]

To bound this last probability, define
\[
K_{\mathrm{tail}}:=K+1,\qquad
\lambda_{\mathrm{act}}:=nb_0qH=\frac{\eta K^2}{\ell}.
\]
The membership and activation pattern has a common law under all the
priors, and each \(V_x\) has distribution
\(\operatorname{Binomial}(n,b_0qH)\). Counting subsets of
\(K_{\mathrm{tail}}\) successful observations gives
\[
E\binom{V_x}{K_{\mathrm{tail}}}
=
\binom{n}{K_{\mathrm{tail}}}(b_0qH)^{K_{\mathrm{tail}}},
\]
with both sides zero when \(K_{\mathrm{tail}}>n\). Since
\(\mathbf1\{V_x>K\}\le\binom{V_x}{K_{\mathrm{tail}}}\),
\[
\Pr(V_x>K)
\le\frac{\lambda_{\mathrm{act}}^{K_{\mathrm{tail}}}}
         {K_{\mathrm{tail}}!}.
\]
The ceiling inequalities imply
\[
\lambda_{\mathrm{act}}
=\frac{\eta K^2}{\ell}
\le2\eta K_{\mathrm{tail}}
\le e^{-31}K_{\mathrm{tail}},
\qquad
K_{\mathrm{tail}}\ge\ell.
\]
Using the \(K_{\mathrm{tail}}\)-th term of the exponential series,
\[
\frac{\lambda_{\mathrm{act}}^{K_{\mathrm{tail}}}}
     {K_{\mathrm{tail}}!}
\le
e^{-31K_{\mathrm{tail}}}
\frac{K_{\mathrm{tail}}^{K_{\mathrm{tail}}}}
     {K_{\mathrm{tail}}!}
\le e^{-30K_{\mathrm{tail}}}
\le e^{-30\ell}.
\]
A union bound, followed by the mass cap and
\(N\le e^\ell\), \(\ell\le e^\ell\), now gives
\[
\widetilde{\mathbb Q}_{\mathrm{act},0}(G_{\mathrm{act}}^c)
\le Je^{-30\ell}
\le \frac{N\ell}{2e^{-32}}e^{-30\ell}
\le e^{32-28\ell}.
\]
Finally, the definition of \(T_{\mathrm{act},\alpha}\) implies
\[
\ell\ge\frac{32-\log(\beta/16)}{28},
\qquad
e^{32-28\ell}\le\frac\beta{16}.
\]
Thus, for every grid index,
\[
\operatorname{TV}
\bigl(\mathbb Q_{\mathrm{act},0},
      \mathbb Q_{\mathrm{act},\iota}\bigr)
\le\frac\beta{16}.
\]
% lean: CausalSmith.Stat.MarRareqLogfrontier.activatedAugmentedMixture_interval_minimax_calibrated

\item \textbf{Coverage of the prior neighborhoods and their length cost.}
Define the target separation, centers, and neighborhood radius by
\[
\Delta_{\mathrm{act}}
:=Jb_0(\mu_{\mathrm{rec},1}-\mu_{\mathrm{rec},0})
\ge\frac{Jb_0}{12}>0,
\]
\[
\theta_{\mathrm{act},\iota}
:=
Jb_0\left[
\left(1-\frac{\iota}{M_{\mathrm{act},\alpha}}\right)
\mu_{\mathrm{rec},0}
+\frac{\iota}{M_{\mathrm{act},\alpha}}\mu_{\mathrm{rec},1}
\right]
=
\theta_{\mathrm{act},0}
+\frac{\iota\Delta_{\mathrm{act}}}{M_{\mathrm{act},\alpha}},
\qquad
w_{\mathrm{act}}
:=\frac{\Delta_{\mathrm{act}}}{8M_{\mathrm{act},\alpha}}.
\]
Under the \(\iota\)-th product prior, the actual target has mean
\(\theta_{\mathrm{act},\iota}\). Independence of the latent coordinates
and \(z_x^{-1}\in[0,1]\) give
\[
\operatorname{Var}\!\left(
b_0\sum_{x=0}^{J-1}z_x^{-1}
\right)\le\frac{Jb_0^2}{4}.
\]
Indeed, a variable in \([0,1]\) has variance at most \(1/4\), and the
variances of these independent summands add. Chebyshev's inequality
therefore gives
\[
\Pr_{\Pi^{(\iota)\otimes d}}\!\left(
\left|\tau(P^{\mathrm{act}}_{\mathbf z})
-\theta_{\mathrm{act},\iota}\right|>w_{\mathrm{act}}
\right)
\le
\frac{Jb_0^2/4}{w_{\mathrm{act}}^2}
\le
\frac{2304M_{\mathrm{act},\alpha}^2}{J}
\le\frac\beta8.
\]

Fix any \(\mathsf I\in\mathcal C_{n,d,q,\alpha}\), and let
\(\nu_{\mathrm{act},\iota}\) denote its returned-interval law when the
sample has law \(\mathbb Q_{\mathrm{act},\iota}\):
\[
\nu_{\mathrm{act},\iota}
:=\mathbb Q_{\mathrm{act},\iota}\mathsf I.
\]
For each grid index, define the interval event
\[
F_{\mathrm{act},\iota}
:=
\left\{
[l,u]:
[l,u]\cap
[\theta_{\mathrm{act},\iota}-w_{\mathrm{act}},
 \theta_{\mathrm{act},\iota}+w_{\mathrm{act}}]
\ne\varnothing
\right\}.
\]
Honesty holds for every fixed completion. Averaging that coverage over
the finite prior and subtracting the target-concentration failure
probability yields
\[
\nu_{\mathrm{act},\iota}(F_{\mathrm{act},\iota})
\ge\beta-\frac\beta8=\frac{7\beta}{8}.
\]
By \cref{obj:lem:randomized-kernel-tv-comparison} and the preceding
likelihood comparison,
\[
\operatorname{TV}
(\nu_{\mathrm{act},0},\nu_{\mathrm{act},\iota})
\le\frac\beta{16},
\qquad
\nu_{\mathrm{act},0}(F_{\mathrm{act},\iota})
\ge\frac{13\beta}{16}.
\]

Define the number of neighborhoods met by a returned interval by
\[
\mathcal N_{\mathrm{hit}}([l,u])
:=
\sum_{\iota=0}^{M_{\mathrm{act},\alpha}}
\mathbf1_{F_{\mathrm{act},\iota}}([l,u]).
\]
Meeting the \(\iota\)-th neighborhood is equivalent to
\(\theta_{\mathrm{act},\iota}\in[l-w_{\mathrm{act}},u+w_{\mathrm{act}}]\).
The centers are spaced by
\(\Delta_{\mathrm{act}}/M_{\mathrm{act},\alpha}\). Thus the distance
between the extreme centers contained in this expanded interval gives
\[
(u-l)+2w_{\mathrm{act}}
\ge
\frac{\Delta_{\mathrm{act}}}{M_{\mathrm{act},\alpha}}
\bigl(\mathcal N_{\mathrm{hit}}([l,u])-1\bigr).
\]
For zero or one contained center, the right side is nonpositive or
zero, so the inequality covers those cases as well. Taking expectations,
\[
\begin{aligned}
E_{\nu_{\mathrm{act},0}}[u-l]
&\ge
\frac{\Delta_{\mathrm{act}}}{M_{\mathrm{act},\alpha}}
\left[
(M_{\mathrm{act},\alpha}+1)\frac{13\beta}{16}-1
\right]
-2w_{\mathrm{act}}\\
&=
\Delta_{\mathrm{act}}
\left[
\frac{13\beta}{16}
+\frac{13\beta}{16M_{\mathrm{act},\alpha}}
-\frac{5}{4M_{\mathrm{act},\alpha}}
\right]\\
&\ge\frac{21\beta}{32}\Delta_{\mathrm{act}},
\end{aligned}
\]
where the last step uses \(1/M_{\mathrm{act},\alpha}\le\beta/8\).
The left side is a prior average of the expected lengths of this same
procedure over laws in \(\mathcal M_{n,d,q}\). Therefore
\[
\frac{21\beta}{32}\Delta_{\mathrm{act}}
\le
\sup_{P\in\mathcal M_{n,d,q}}E_P[\operatorname{length}(\mathsf I)].
\]
The constant procedure returning \([-1,1]\) is honest, so the procedure
class is nonempty. Taking its defining infimum gives
\[
\frac{21\beta}{32}\Delta_{\mathrm{act}}
\le\mathfrak L_{n,d,q,\alpha}.
\]

To express the separation in terms of \(d\), split according to the
minimum defining \(J\). If
\(d-1\le\lfloor(2b_0)^{-1}\rfloor\), then \(J=d-1\), and
\(J\ge D_{\mathrm{act},\alpha}\ge1\) implies \(d\ge2\) and
\(J\ge d/2\). Otherwise,
\(J=\lfloor(2b_0)^{-1}\rfloor\), and
\[
Jb_0\ge\frac12-b_0\ge\frac14.
\]
In both cases,
\[
Jb_0\ge\min\left\{\frac{db_0}{2},\frac14\right\}.
\]
Since \(\eta\le1/2\), the two inequalities
\[
\frac{\eta}{24}s_{\mathrm{many}}
\le\frac{db_0}{24},
\qquad
\frac{\eta}{24}s_{\mathrm{many}}
\le\frac1{48}
\]
give
\[
\frac{\eta}{24}s_{\mathrm{many}}
\le
\frac1{12}\min\left\{\frac{db_0}{2},\frac14\right\}
\le\frac{Jb_0}{12}
\le\Delta_{\mathrm{act}}.
\]
Consequently,
\[
c_{\mathrm{act},\alpha}s_{\mathrm{many}}
\le\mathfrak L_{n,d,q,\alpha}
\qquad
(N\ge T_{\mathrm{act},\alpha},\ J\ge D_{\mathrm{act},\alpha}).
\]
% lean: CausalSmith.Stat.MarRareqLogfrontier.interval_activated_large_regime

\item \textbf{An explicit one-cell floor.}
Use the family \(P_u\) supplied by
\cref{obj:lem:one-cell-testing-family}. We record the calibration needed
here:
\[
c_{\mathrm{par},\alpha}:=\frac{3\beta^2}{512}>0,
\qquad
c_{\mathrm{par},\alpha}s_{\mathrm{par}}
\le\mathfrak L_{n,d,q,\alpha}.
\]
To obtain this explicit constant, define
\[
\eta_{\mathrm{cell}}:=\frac\beta8,\qquad
w_{\mathrm{cell}}
:=\min\left\{\frac14,\frac{\eta_{\mathrm{cell}}}{4\sqrt N}\right\}.
\]
Choose an integer \(M_{\mathrm{cell},\alpha}\ge8/\beta\), and set
\[
\delta_{\mathrm{cell}}
:=\frac{w_{\mathrm{cell}}}{M_{\mathrm{cell},\alpha}}>0,
\qquad
u_{\mathrm{cell},\iota}
:=\frac12+\iota\delta_{\mathrm{cell}}
\quad(0\le\iota\le M_{\mathrm{cell},\alpha}).
\]
These grid points are distinct and belong to \([1/4,3/4]\), and
\[
2Nw_{\mathrm{cell}}^2\le\frac{\eta_{\mathrm{cell}}^2}{8},
\qquad
w_{\mathrm{cell}}\ge
\frac{\eta_{\mathrm{cell}}}{4}s_{\mathrm{par}}.
\]
Indeed, \(N>0\) and \(0<\eta_{\mathrm{cell}}=\beta/8<1\), so
\(w_{\mathrm{cell}}>0\). The defining minimum gives
\[
w_{\mathrm{cell}}\le\frac{\eta_{\mathrm{cell}}}{4\sqrt N},
\qquad
2Nw_{\mathrm{cell}}^2
\le 2N\left(\frac{\eta_{\mathrm{cell}}}{4\sqrt N}\right)^2
=\frac{\eta_{\mathrm{cell}}^2}{8},
\]
where the equality uses \((\sqrt N)^2=N\). For the lower bound, if
\(1/4\le\eta_{\mathrm{cell}}/(4\sqrt N)\), then
\(w_{\mathrm{cell}}=1/4\), and \(s_{\mathrm{par}}\le1\) gives
\[
\frac{\eta_{\mathrm{cell}}}{4}s_{\mathrm{par}}
\le\frac{\eta_{\mathrm{cell}}}{4}
\le\frac14=w_{\mathrm{cell}}.
\]
Otherwise, \(w_{\mathrm{cell}}=\eta_{\mathrm{cell}}/(4\sqrt N)\),
and \(s_{\mathrm{par}}\le N^{-1/2}\) gives
\[
\frac{\eta_{\mathrm{cell}}}{4}s_{\mathrm{par}}
\le\frac{\eta_{\mathrm{cell}}}{4\sqrt N}
=w_{\mathrm{cell}}.
\]

Write the one-record and sample laws of this family as
\[
\mathbb Q_{\mathrm{cell},u}^{(1)}:=\mathcal L_{P_u}(O_1),
\qquad
\mathbb Q_{\mathrm{cell},u}^{(n)}
:=(\mathbb Q_{\mathrm{cell},u}^{(1)})^{\otimes n}.
\]
In this family, the arrival floor and \(P_u(R=1)=q\), together with
balanced assignment, force arrival probability \(q\) in each treatment
arm. Conditional independence of arrival and outcome then gives the
two treated arrived-outcome atom probabilities \(qu/2\) and
\(q(1-u)/2\); all other observed atom probabilities are constant in \(u\).
Consequently, defining the one-record chi-squared divergence by its
finite sum,
\[
\mathcal X_{\mathrm{cell},\iota}^{(1)}
:=
\sum_{\substack{o:\\
\mathbb Q_{\mathrm{cell},1/2}^{(1)}(\{o\})>0}}
\frac{
\left[
\mathbb Q_{\mathrm{cell},u_{\mathrm{cell},\iota}}^{(1)}(\{o\})
-\mathbb Q_{\mathrm{cell},1/2}^{(1)}(\{o\})
\right]^2}
{\mathbb Q_{\mathrm{cell},1/2}^{(1)}(\{o\})}
=
2q(\iota\delta_{\mathrm{cell}})^2.
\]
This computation also covers \(q=1\), since the vanished nonarrival
atoms are omitted from the sum. Define
\(\mathcal X_{\mathrm{cell},\iota}^{(n)}\) by the analogous sum for the
sample laws, and put
\[
t_{\mathrm{cell},\iota}
:=2N(\iota\delta_{\mathrm{cell}})^2
\le\frac{\eta_{\mathrm{cell}}^2}{8}.
\]
Independence of the observations gives
\[
1+\mathcal X_{\mathrm{cell},\iota}^{(n)}
=
\bigl(1+\mathcal X_{\mathrm{cell},\iota}^{(1)}\bigr)^n
\le e^{t_{\mathrm{cell},\iota}}.
\]
For \(0\le t_{\mathrm{cell},\iota}<1\), comparison of the exponential
and geometric series gives
\[
e^{t_{\mathrm{cell},\iota}}
\le\frac1{1-t_{\mathrm{cell},\iota}},
\qquad
\frac{t_{\mathrm{cell},\iota}}{1-t_{\mathrm{cell},\iota}}
\le\eta_{\mathrm{cell}}^2,
\]
using \(t_{\mathrm{cell},\iota}\le\eta_{\mathrm{cell}}^2/8\) and
\(\eta_{\mathrm{cell}}\le1\). Thus Cauchy--Schwarz applied to the
finite likelihood-ratio sum yields
\[
\operatorname{TV}
\bigl(\mathbb Q_{\mathrm{cell},1/2}^{(n)},
      \mathbb Q_{\mathrm{cell},u_{\mathrm{cell},\iota}}^{(n)}\bigr)
\le\frac12\sqrt{\mathcal X_{\mathrm{cell},\iota}^{(n)}}
\le\eta_{\mathrm{cell}}.
\]

For any honest procedure \(\mathsf I\), coverage of each grid point is
at least \(\beta\) under its corresponding sample law. By
\cref{obj:lem:randomized-kernel-tv-comparison}, under the reference law
\(P_{1/2}\) each point is contained in the returned interval with
probability at least
\(\beta-\eta_{\mathrm{cell}}=7\beta/8\). A connected interval
containing a given number of grid points has length at least their
spacing times that number minus one. Summing the inclusion
probabilities therefore gives
\[
\begin{aligned}
E_{P_{1/2}}[\operatorname{length}(\mathsf I)]
&\ge
\delta_{\mathrm{cell}}
\left[(M_{\mathrm{cell},\alpha}+1)\frac{7\beta}{8}-1\right]\\
&\ge\frac{3\beta}{4}w_{\mathrm{cell}},
\end{aligned}
\]
because \(1/M_{\mathrm{cell},\alpha}\le\beta/8\).
All interval risks lie in \([0,2]\), and the constant interval
\([-1,1]\) is honest. Hence taking the supremum over models and the
infimum over honest procedures preserves this lower bound. Finally,
\[
c_{\mathrm{par},\alpha}s_{\mathrm{par}}
\le
\frac{3\beta}{4}\frac{\eta_{\mathrm{cell}}}{4}s_{\mathrm{par}}
\le
\frac{3\beta}{4}w_{\mathrm{cell}}
\le\mathfrak L_{n,d,q,\alpha},
\]
which proves the asserted explicit floor.
% lean: CausalSmith.Stat.MarRareqLogfrontier.oneCell_interval_minimax_floor

\item \textbf{Extension of the many-cell scale to all regimes.}
Define
\[
c_{\mathrm{many},\alpha}
:=
\min\left\{
c_{\mathrm{act},\alpha},
\ c_{\mathrm{par},\alpha}
       \min\{1,T_{\mathrm{act},\alpha}^{-1/2}\},
\ \frac{c_{\mathrm{par},\alpha}}{D_{\mathrm{act},\alpha}}
\right\}>0.
\]
We claim
\[
c_{\mathrm{many},\alpha}s_{\mathrm{many}}
\le\mathfrak L_{n,d,q,\alpha}
\]
throughout the parameter range.

First suppose \(N\ge T_{\mathrm{act},\alpha}\). If
\(J\ge D_{\mathrm{act},\alpha}\), the activated bound gives
\[
c_{\mathrm{many},\alpha}s_{\mathrm{many}}
\le c_{\mathrm{act},\alpha}s_{\mathrm{many}}
\le\mathfrak L_{n,d,q,\alpha}.
\]
If \(J<D_{\mathrm{act},\alpha}\), the capacity inequality from \cref{proof:connected-interval-calibration}
implies
\[
\left\lfloor\frac1{2b_0}\right\rfloor
\ge D_{\mathrm{act},\alpha}.
\]
Thus the minimum defining \(J\) forces
\(d-1<D_{\mathrm{act},\alpha}\), and hence
\(d\le D_{\mathrm{act},\alpha}\). Since \(N\ge1\) and \(\ell\ge1\),
\[
s_{\mathrm{many}}
\le\frac{d}{N\ell}
\le\frac{D_{\mathrm{act},\alpha}}{\sqrt N}
=
D_{\mathrm{act},\alpha}s_{\mathrm{par}}.
\]
It follows that
\[
c_{\mathrm{many},\alpha}s_{\mathrm{many}}
\le
c_{\mathrm{many},\alpha}D_{\mathrm{act},\alpha}s_{\mathrm{par}}
\le c_{\mathrm{par},\alpha}s_{\mathrm{par}}
\le\mathfrak L_{n,d,q,\alpha}.
\]
This case includes \(d=1\) and \(J=0\).

Finally, if \(N<T_{\mathrm{act},\alpha}\), monotonicity of the square
root and reciprocal gives
\[
\min\{1,T_{\mathrm{act},\alpha}^{-1/2}\}
\le s_{\mathrm{par}}.
\]
Using \(s_{\mathrm{many}}\le1\), we obtain
\[
c_{\mathrm{many},\alpha}s_{\mathrm{many}}
\le c_{\mathrm{many},\alpha}
\le
c_{\mathrm{par},\alpha}
\min\{1,T_{\mathrm{act},\alpha}^{-1/2}\}
\le\mathfrak L_{n,d,q,\alpha}.
\]
% lean: CausalSmith.Stat.MarRareqLogfrontier.interval_many_scale_of_large_regime

\item \label{proof:connected-interval-lower-scales}\textbf{Combination of the lower scales.}
The identities
\[
\sqrt{\min\{1,N^{-1}\}}=s_{\mathrm{par}},
\qquad
\sqrt{\min\left\{1,\left(\frac{d}{N\ell}\right)^2\right\}}
=s_{\mathrm{many}}
\]
follow by splitting at \(N^{-1}=1\) and \(d/(N\ell)=1\), respectively.
Define
\[
c_{\mathrm{both},\alpha}
:=\min\{c_{\mathrm{par},\alpha},c_{\mathrm{many},\alpha}\}>0.
\]
The preceding bounds give
\[
c_{\mathrm{both},\alpha}
\max\{s_{\mathrm{par}},s_{\mathrm{many}}\}
\le\mathfrak L_{n,d,q,\alpha}.
\]
Moreover,
\[
r_{n,d,q}
\le2\max\{s_{\mathrm{par}}^2,s_{\mathrm{many}}^2\}.
\]
Indeed, when both untruncated terms \(N^{-1}\) and
\((d/(N\ell))^2\) are at most one, their sum is at most twice their
maximum. If either exceeds one, the maximum of the truncated terms
equals one, while \(r_{n,d,q}\le1\). Since the maximum is nonnegative,
\[
\sqrt{r_{n,d,q}}
\le2\max\{s_{\mathrm{par}},s_{\mathrm{many}}\}.
\]
Therefore, defining
\[
c_{\mathrm{low},\alpha}:=\frac{c_{\mathrm{both},\alpha}}2>0,
\]
we obtain
\[
c_{\mathrm{low},\alpha}\sqrt{r_{n,d,q}}
\le\mathfrak L_{n,d,q,\alpha}.
\]
% lean: CausalSmith.Stat.MarRareqLogfrontier.frontier_interval_lower_of_two_scales

\item \textbf{Honesty and the minimax comparison.}
Set
\[
\underline c_\alpha
:=\min\{c_{\mathrm{low},\alpha},C_\alpha\},
\qquad
\overline C_\alpha:=C_\alpha.
\]
Then \(0<\underline c_\alpha\le\overline C_\alpha<\infty\), and
\cref{proof:connected-interval-lower-scales} supplies the required lower inequality.

Fix \(P\in\mathcal M_{n,d,q}\). Define the squared error by
\[
V_{\mathrm{env}}
:=
\bigl(\widehat\tau^{\mathrm{mix}}_{n,d,q}(\mathbf O_n)-\tau(P)\bigr)^2.
\]
The risk bound implies
\[
0\le E_P[V_{\mathrm{env}}]\le\mathfrak U_{n,d,q}.
\]
If \(\mathfrak U_{n,d,q}=0\), nonnegativity gives
\(V_{\mathrm{env}}=0\) almost surely, so the interval contains
\(\tau(P)\) almost surely. If \(\mathfrak U_{n,d,q}>0\), define
\[
t_{\mathrm{env}}:=\frac{\mathfrak U_{n,d,q}}{\alpha}>0.
\]
Integrating the pointwise inequality
\(V_{\mathrm{env}}\ge
t_{\mathrm{env}}\mathbf1\{V_{\mathrm{env}}\ge t_{\mathrm{env}}\}\)
gives
\[
t_{\mathrm{env}}\Pr_P(V_{\mathrm{env}}\ge t_{\mathrm{env}})
\le E_P[V_{\mathrm{env}}]\le\mathfrak U_{n,d,q},
\qquad
\Pr_P(V_{\mathrm{env}}\ge t_{\mathrm{env}})\le\alpha.
\]
Because \(\tau(P)\in[-1,1]\), squared error at most
\(t_{\mathrm{env}}\) implies inclusion in the clipped interval.
Consequently,
\[
\{\tau(P)\notin I^{\mathrm{mix}}_\alpha\}
\subseteq\{V_{\mathrm{env}}\ge t_{\mathrm{env}}\},
\qquad
\Pr_P\{\tau(P)\in I^{\mathrm{mix}}_\alpha\}\ge1-\alpha.
\]

The mixed-count estimator takes values in \([-1,1]\) by
\cref{obj:def:mixed-count-estimator}. Its envelope interval therefore
has ordered endpoints in \([-1,1]\). The finite observed-sample space
makes its endpoint map measurable, and its Dirac kernel is an honest
member of \(\mathcal C_{n,d,q,\alpha}\). For this kernel, joint
sample-and-kernel expected length equals
\(E_P[\operatorname{length}(I^{\mathrm{mix}}_\alpha)]\).
The definition in \cref{obj:def:interval-length-risk} hence gives
\[
\mathfrak L_{n,d,q,\alpha}
\le
\sup_{P\in\mathcal M_{n,d,q}}
E_P[\operatorname{length}(I^{\mathrm{mix}}_\alpha)].
\]
Combining this comparison with \cref{proof:connected-interval-upper,proof:connected-interval-lower-scales} proves
\[
\underline c_\alpha\sqrt{r_{n,d,q}}
\le\mathfrak L_{n,d,q,\alpha}
\le
\sup_{P\in\mathcal M_{n,d,q}}
E_P[\operatorname{length}(I^{\mathrm{mix}}_\alpha)]
\le\overline C_\alpha\sqrt{r_{n,d,q}},
\]
together with the asserted coverage.
% lean: CausalSmith.Stat.MarRareqLogfrontier.connected_interval_frontier
\end{enumerate}
\end{proof}
